\clearpage\section{MOS transistor schematic}
These editable vector schematics connect equal node names within a circuit. Primitive-cell names are local pins; the hierarchy defines their connections. B denotes the body terminal; 0 is global ground. Each panel identifies its netlist and scope; measurement sources are connected as described.
Complete six-MOS baseline differential amplifier from \texttt{01\_basic\_differential.sp}.
\par
\begin{figure}[!ht]
\centering
\begin{circuitikz}[font=\small]
\node[pmos,bulk] (M3) at (0,4.7) {};
\node[anchor=west,font=\scriptsize] at ($(M3.center)+(.55,.40)$) {M3};
\draw (M3.bulk) -- ++(.30,0) node[right,font=\scriptsize] {B=vdd};
\node[pmos,bulk] (M4) at (4.2,4.7) {};
\node[anchor=west,font=\scriptsize] at ($(M4.center)+(.55,.40)$) {M4};
\draw (M4.bulk) -- ++(.30,0) node[right,font=\scriptsize] {B=vdd};
\node[nmos,bulk] (M1) at (0,2.4) {};
\node[anchor=west,font=\scriptsize] at ($(M1.center)+(.55,.40)$) {M1};
\draw (M1.bulk) -- ++(.30,0) node[right,font=\scriptsize] {B=vss};
\draw (M1.G) -- ++(-.45,0) node[left,font=\small] {\texttt{ip}};
\node[nmos,bulk] (M2) at (4.2,2.4) {};
\node[anchor=west,font=\scriptsize] at ($(M2.center)+(.55,.40)$) {M2};
\draw (M2.bulk) -- ++(.30,0) node[right,font=\scriptsize] {B=vss};
\draw (M2.G) -- ++(-.45,0) node[left,font=\small] {\texttt{im}};
\draw (M3.S) -- ++(0,0.5) node[above,font=\small] {\texttt{vdd}};
\draw (M4.S) -- ++(0,0.5) node[above,font=\small] {\texttt{vdd}};
\draw (M3.D) -- (M1.D) node[midway,above right,font=\small] {\texttt{left}};
\draw (M4.D) -- (M2.D) node[midway,above right,font=\small] {\texttt{out}};
\draw (M3.G) -- (-1,4.7) -- (-1,3.55) -- (M3.D |- 0,3.55) node[circ]{}; \draw (M4.G) -- (2.9,4.7) -- (2.9,3.55) -- (M3.D |- 0,3.55);
\node[nmos,bulk] (MTAIL) at (2.1,-0.3) {};
\node[anchor=west,font=\scriptsize] at ($(MTAIL.center)+(.55,.40)$) {MTAIL};
\draw (MTAIL.bulk) -- ++(.30,0) node[right,font=\scriptsize] {B=vss};
\draw (MTAIL.G) -- ++(-.45,0) node[left,font=\small] {\texttt{bn}};
\node[nmos,bulk] (MREF) at (5.6,-0.3) {};
\node[anchor=west,font=\scriptsize] at ($(MREF.center)+(.55,.40)$) {MREF};
\draw (MREF.bulk) -- ++(.30,0) node[right,font=\scriptsize] {B=vss};
\draw (MREF.G) -- ++(-.45,0) node[left,font=\small] {\texttt{bn}};
\draw (M1.S) -- (M1.S |- 0,1.2) -- (M2.S |- 0,1.2) -- (M2.S); \draw (MTAIL.D) -- (MTAIL.D |- 0,1.2) node[circ,label=above:\texttt{tail}]{};
\draw (MTAIL.S) -- ++(0,-0.4) node[below,font=\small] {\texttt{vss}};
\draw (MREF.D) -- ++(0,0.5) node[above,font=\small] {\texttt{bn}};
\draw (MREF.S) -- ++(0,-0.4) node[below,font=\small] {\texttt{vss}};
\draw (8.2,5.4) node[above] {\texttt{vdd}} to[I,l={IRE}] ++(0,-1.7) node[below] {\texttt{bn}};
\draw (8.2,1.5) node[above] {\texttt{out}} to[C,l={CL}] ++(0,-1.7) node[below] {\texttt{0}};
\end{circuitikz}
\caption{NMOS differential pair, PMOS mirror load and NMOS tail-current mirror.}
\end{figure}
\par
All NMOS bodies connect to vss, including M1/M2 whose sources are tail. All PMOS sources/bodies connect to vdd. At scale=1 the widths are MREF=85, MTAIL=870, M1=M2=425 and M3=M4=1350 $\mu$m; all lengths are 0.3 $\mu$m. The scale parameter multiplies both W and L of M1--M4 only. IRE is 50 $\mu$A, CL is 10 pF. The input sources impose $ip=cmnode+diff/2$ and $im=cmnode-diff/2$. vdd=1 V and vss=0 V at DC; CL returns to global 0, which is distinct from an AC-perturbed vss in rejection tests. Cascode-load, full-cascode and current-amplifier decks are separate variants; the figure covers the named baseline deck.
